% --- Archon physics formalization source begin ---
% archon:physics
% archon:covers PhyXMiniProblems/problem_phyx_mini_0204.lean
% archon:source-report reports/phyx_mini/problem_phyx_mini_0204.source.json

\chapter{Physics problem phyx\_mini\_0204}
\label{ch:PhyXMiniProblems_problem_phyx_mini_0204}

\paragraph{Problem source.}
\#\# Physical scenario

One end of a horizontal string is attached to a small-amplitude mechanical 60.0\textbackslash{},\textbackslash{}mathrm\{Hz\} oscillator. The string's mass per unit length is \textbackslash{}( 3.5 \textbackslash{}times 10\textasciicircum{}\{-4\}\textbackslash{},\textbackslash{}mathrm\{kg/m\} \textbackslash{}). The string passes over a pulley, a distance \textbackslash{}( \textbackslash{}ell = 1.50\textbackslash{},\textbackslash{}mathrm\{m\} \textbackslash{}) away, and weights are hung from this end, figure.Assume the string at the oscillator is a node, which is nearly true.

\#\# Question

What mass \textbackslash{}( m \textbackslash{}) must be hung from this end of the string to produce one loop?

\#\# Answer choices

- A: A: 1.8kg

- B: B: 1.6kg

- C: C: 1.4kg

- D: D: 1.2kg

\#\# Figure caption (auxiliary; use the image as primary evidence)

The image depicts a physical setup involving an oscillator and a pulley system connected to a mass. Here's a detailed description:

- On the left side, there is a rectangular block labeled "Oscillator." It is connected to a horizontal rope or string.
- The rope extends horizontally for 1.50 meters, as indicated by a double-headed arrow and text above the rope.
- The rope then passes over a pulley positioned at the right end of the horizontal section.
- The pulley redirects the rope vertically downward.
- At the end of the vertical section of the rope, there is a block labeled with the letter "m," representing the mass.
- The pulley system allows the mass to be lifted or suspended by the oscillator through the rope.

The image illustrates the mechanical relationship between the oscillator, the pulley, and the mass.

\#\# Dataset metadata

Resonance and Harmonics; Physical Model Grounding Reasoning, Multi-Formula Reasoning

\paragraph{Recorded answer/context.}
D: D: 1.2kg

\paragraph{Figure/image path.}
/root/proposal\_for\_physic/science-mango/phyx\_mini\_run/phyx\_data/test\_image/204.png

\paragraph{Formalization target.}
create a compiling Lean file with sorry bodies at `PhyXMiniProblems/problem\_phyx\_mini\_0204.lean`. The Lean declarations must preserve the physical quantities, dimensions or dimensional roles, figure labels, governing-law hypotheses, and final relation expressed by this problem.
Use Mathlib/Physlib names found through LeanExplore where available. If a physics API is missing, introduce faithful local abstractions rather than scalar placeholder aliases.

\begin{theorem}[Physics formalization target]
\label{thm:physics:phyx_mini_0204:target}
\lean{PhyXMiniProblems.ProblemPhyXMini0204.hangingMass_matches_recordedAnswerD}
\uses{def:physics:phyx-mini-0204:phyxminiproblems-problemphyxmini0204-stringpulleysetup, def:physics:phyx-mini-0204:phyxminiproblems-problemphyxmini0204-matchesproblemandprimaryfigure, def:physics:phyx-mini-0204:phyxminiproblems-problemphyxmini0204-hasnodeboundaryconditions, def:physics:phyx-mini-0204:phyxminiproblems-problemphyxmini0204-isrequestedoneloopstate, def:physics:phyx-mini-0204:phyxminiproblems-problemphyxmini0204-usesstandardgravity, def:physics:phyx-mini-0204:phyxminiproblems-problemphyxmini0204-haspositivephysicalparameters, def:physics:phyx-mini-0204:phyxminiproblems-problemphyxmini0204-satisfiesnodetonodestandingwavelaw, def:physics:phyx-mini-0204:phyxminiproblems-problemphyxmini0204-satisfieswavespeedlaw, def:physics:phyx-mini-0204:phyxminiproblems-problemphyxmini0204-satisfiestautstringtensionlaw, def:physics:phyx-mini-0204:phyxminiproblems-problemphyxmini0204-satisfiesidealpulleyandstaticloadlaw, def:physics:phyx-mini-0204:phyxminiproblems-problemphyxmini0204-recordedanswerchoice, def:physics:phyx-mini-0204:phyxminiproblems-problemphyxmini0204-matchesanswerchoice}
The assigned autoformalize agent should translate this physics problem into Lean declarations in the covered file, with theorem and lemma proof bodies written as `by sorry`.
\end{theorem}
\begin{proof}
This is an autoformalization task, not a proof task. Produce statements that can later be proved by the physics prover without weakening the physical model.
\end{proof}

% --- Archon declaration topology begin ---
\section{Declaration topology}
\begin{definition}[Length Quantity]
  \label{def:physics:phyx-mini-0204:phyxminiproblems-problemphyxmini0204-lengthquantity}
  \lean{PhyXMiniProblems.ProblemPhyXMini0204.LengthQuantity}
  A physical length
\end{definition}
\begin{proof}
  This is the declared model definition of Length Quantity.
\end{proof}
\begin{definition}[Mass Quantity]
  \label{def:physics:phyx-mini-0204:phyxminiproblems-problemphyxmini0204-massquantity}
  \lean{PhyXMiniProblems.ProblemPhyXMini0204.MassQuantity}
  A physical mass
\end{definition}
\begin{proof}
  This is the declared model definition of Mass Quantity.
\end{proof}
\begin{definition}[Frequency Quantity]
  \label{def:physics:phyx-mini-0204:phyxminiproblems-problemphyxmini0204-frequencyquantity}
  \lean{PhyXMiniProblems.ProblemPhyXMini0204.FrequencyQuantity}
  A physical frequency, carrying the inverse-time dimension
\end{definition}
\begin{proof}
  This is the declared model definition of Frequency Quantity.
\end{proof}
\begin{definition}[Linear Mass Density Quantity]
  \label{def:physics:phyx-mini-0204:phyxminiproblems-problemphyxmini0204-linearmassdensityquantity}
  \lean{PhyXMiniProblems.ProblemPhyXMini0204.LinearMassDensityQuantity}
  A physical mass per unit length
\end{definition}
\begin{proof}
  This is the declared model definition of Linear Mass Density Quantity.
\end{proof}
\begin{definition}[Speed Quantity]
  \label{def:physics:phyx-mini-0204:phyxminiproblems-problemphyxmini0204-speedquantity}
  \lean{PhyXMiniProblems.ProblemPhyXMini0204.SpeedQuantity}
  A physical propagation speed
\end{definition}
\begin{proof}
  This is the declared model definition of Speed Quantity.
\end{proof}
\begin{definition}[Force Quantity]
  \label{def:physics:phyx-mini-0204:phyxminiproblems-problemphyxmini0204-forcequantity}
  \lean{PhyXMiniProblems.ProblemPhyXMini0204.ForceQuantity}
  A physical force, used for the two string-tension magnitudes
\end{definition}
\begin{proof}
  This is the declared model definition of Force Quantity.
\end{proof}
\begin{definition}[Acceleration Quantity]
  \label{def:physics:phyx-mini-0204:phyxminiproblems-problemphyxmini0204-accelerationquantity}
  \lean{PhyXMiniProblems.ProblemPhyXMini0204.AccelerationQuantity}
  A physical acceleration magnitude
\end{definition}
\begin{proof}
  This is the declared model definition of Acceleration Quantity.
\end{proof}
\begin{definition}[quantity Readout]
  \label{def:physics:phyx-mini-0204:phyxminiproblems-problemphyxmini0204-quantityreadout}
  \lean{PhyXMiniProblems.ProblemPhyXMini0204.quantityReadout}
  Read a dimension-tagged nonnegative physical quantity in chosen units
\end{definition}
\begin{proof}
  This is the declared model definition of quantity Readout.
\end{proof}
\begin{definition}[length In Meters]
  \label{def:physics:phyx-mini-0204:phyxminiproblems-problemphyxmini0204-lengthinmeters}
  \lean{PhyXMiniProblems.ProblemPhyXMini0204.lengthInMeters}
  \uses{def:physics:phyx-mini-0204:phyxminiproblems-problemphyxmini0204-lengthquantity, def:physics:phyx-mini-0204:phyxminiproblems-problemphyxmini0204-quantityreadout}
  Metre readout of a physical length
\end{definition}
\begin{proof}
  This is the declared model definition of length In Meters.
\end{proof}
\begin{definition}[mass In Kilograms]
  \label{def:physics:phyx-mini-0204:phyxminiproblems-problemphyxmini0204-massinkilograms}
  \lean{PhyXMiniProblems.ProblemPhyXMini0204.massInKilograms}
  \uses{def:physics:phyx-mini-0204:phyxminiproblems-problemphyxmini0204-massquantity, def:physics:phyx-mini-0204:phyxminiproblems-problemphyxmini0204-quantityreadout}
  Kilogram readout of a physical mass
\end{definition}
\begin{proof}
  This is the declared model definition of mass In Kilograms.
\end{proof}
\begin{definition}[frequency In Hertz]
  \label{def:physics:phyx-mini-0204:phyxminiproblems-problemphyxmini0204-frequencyinhertz}
  \lean{PhyXMiniProblems.ProblemPhyXMini0204.frequencyInHertz}
  \uses{def:physics:phyx-mini-0204:phyxminiproblems-problemphyxmini0204-frequencyquantity, def:physics:phyx-mini-0204:phyxminiproblems-problemphyxmini0204-quantityreadout}
  Hertz readout of a physical frequency
\end{definition}
\begin{proof}
  This is the declared model definition of frequency In Hertz.
\end{proof}
\begin{definition}[linear Mass Density In Kilograms Per Meter]
  \label{def:physics:phyx-mini-0204:phyxminiproblems-problemphyxmini0204-linearmassdensityinkilogramspermeter}
  \lean{PhyXMiniProblems.ProblemPhyXMini0204.linearMassDensityInKilogramsPerMeter}
  \uses{def:physics:phyx-mini-0204:phyxminiproblems-problemphyxmini0204-linearmassdensityquantity, def:physics:phyx-mini-0204:phyxminiproblems-problemphyxmini0204-quantityreadout}
  Kilograms-per-metre readout of a linear mass density
\end{definition}
\begin{proof}
  This is the declared model definition of linear Mass Density In Kilograms Per Meter.
\end{proof}
\begin{definition}[speed In Meters Per Second]
  \label{def:physics:phyx-mini-0204:phyxminiproblems-problemphyxmini0204-speedinmeterspersecond}
  \lean{PhyXMiniProblems.ProblemPhyXMini0204.speedInMetersPerSecond}
  \uses{def:physics:phyx-mini-0204:phyxminiproblems-problemphyxmini0204-speedquantity, def:physics:phyx-mini-0204:phyxminiproblems-problemphyxmini0204-quantityreadout}
  Metres-per-second readout of a physical speed
\end{definition}
\begin{proof}
  This is the declared model definition of speed In Meters Per Second.
\end{proof}
\begin{definition}[force In Newtons]
  \label{def:physics:phyx-mini-0204:phyxminiproblems-problemphyxmini0204-forceinnewtons}
  \lean{PhyXMiniProblems.ProblemPhyXMini0204.forceInNewtons}
  \uses{def:physics:phyx-mini-0204:phyxminiproblems-problemphyxmini0204-forcequantity, def:physics:phyx-mini-0204:phyxminiproblems-problemphyxmini0204-quantityreadout}
  Newton readout of a force in coherent SI base units
\end{definition}
\begin{proof}
  This is the declared model definition of force In Newtons.
\end{proof}
\begin{definition}[acceleration In Meters Per Second Squared]
  \label{def:physics:phyx-mini-0204:phyxminiproblems-problemphyxmini0204-accelerationinmeterspersecondsquared}
  \lean{PhyXMiniProblems.ProblemPhyXMini0204.accelerationInMetersPerSecondSquared}
  \uses{def:physics:phyx-mini-0204:phyxminiproblems-problemphyxmini0204-accelerationquantity, def:physics:phyx-mini-0204:phyxminiproblems-problemphyxmini0204-quantityreadout}
  Metres-per-second-squared readout of an acceleration magnitude
\end{definition}
\begin{proof}
  This is the declared model definition of acceleration In Meters Per Second Squared.
\end{proof}
\begin{definition}[Figure Component]
  \label{def:physics:phyx-mini-0204:phyxminiproblems-problemphyxmini0204-figurecomponent}
  \lean{PhyXMiniProblems.ProblemPhyXMini0204.FigureComponent}
  Individually visible components of the supplied apparatus figure
\end{definition}
\begin{proof}
  This is the declared model definition of Figure Component.
\end{proof}
\begin{definition}[Horizontal Span Endpoint]
  \label{def:physics:phyx-mini-0204:phyxminiproblems-problemphyxmini0204-horizontalspanendpoint}
  \lean{PhyXMiniProblems.ProblemPhyXMini0204.HorizontalSpanEndpoint}
  The two ends of the horizontally vibrating string span
\end{definition}
\begin{proof}
  This is the declared model definition of Horizontal Span Endpoint.
\end{proof}
\begin{definition}[String Route]
  \label{def:physics:phyx-mini-0204:phyxminiproblems-problemphyxmini0204-stringroute}
  \lean{PhyXMiniProblems.ProblemPhyXMini0204.StringRoute}
  The route of the string visible in the primary figure
\end{definition}
\begin{proof}
  This is the declared model definition of String Route.
\end{proof}
\begin{definition}[Oscillator Regime]
  \label{def:physics:phyx-mini-0204:phyxminiproblems-problemphyxmini0204-oscillatorregime}
  \lean{PhyXMiniProblems.ProblemPhyXMini0204.OscillatorRegime}
  The stated operating regime of the oscillator
\end{definition}
\begin{proof}
  This is the declared model definition of Oscillator Regime.
\end{proof}
\begin{definition}[String Pulley Setup]
  \label{def:physics:phyx-mini-0204:phyxminiproblems-problemphyxmini0204-stringpulleysetup}
  \lean{PhyXMiniProblems.ProblemPhyXMini0204.StringPulleySetup}
  \uses{def:physics:phyx-mini-0204:phyxminiproblems-problemphyxmini0204-lengthquantity, def:physics:phyx-mini-0204:phyxminiproblems-problemphyxmini0204-massquantity, def:physics:phyx-mini-0204:phyxminiproblems-problemphyxmini0204-frequencyquantity, def:physics:phyx-mini-0204:phyxminiproblems-problemphyxmini0204-linearmassdensityquantity, def:physics:phyx-mini-0204:phyxminiproblems-problemphyxmini0204-speedquantity, def:physics:phyx-mini-0204:phyxminiproblems-problemphyxmini0204-forcequantity, def:physics:phyx-mini-0204:phyxminiproblems-problemphyxmini0204-accelerationquantity, def:physics:phyx-mini-0204:phyxminiproblems-problemphyxmini0204-figurecomponent, def:physics:phyx-mini-0204:phyxminiproblems-problemphyxmini0204-horizontalspanendpoint, def:physics:phyx-mini-0204:phyxminiproblems-problemphyxmini0204-stringroute, def:physics:phyx-mini-0204:phyxminiproblems-problemphyxmini0204-oscillatorregime}
  The physical quantities and categorical data of the oscillator--string--pulley apparatus. loopCount, wavelength, speed, both tensions, and hanging mass are unknown state variables; this structure assigns none of them the requested answer value
\end{definition}
\begin{proof}
  This is the declared model definition of String Pulley Setup.
\end{proof}
\begin{definition}[Matches Problem And Primary Figure]
  \label{def:physics:phyx-mini-0204:phyxminiproblems-problemphyxmini0204-matchesproblemandprimaryfigure}
  \lean{PhyXMiniProblems.ProblemPhyXMini0204.MatchesProblemAndPrimaryFigure}
  \uses{def:physics:phyx-mini-0204:phyxminiproblems-problemphyxmini0204-lengthinmeters, def:physics:phyx-mini-0204:phyxminiproblems-problemphyxmini0204-frequencyinhertz, def:physics:phyx-mini-0204:phyxminiproblems-problemphyxmini0204-linearmassdensityinkilogramspermeter, def:physics:phyx-mini-0204:phyxminiproblems-problemphyxmini0204-stringpulleysetup}
  The categorical and numerical data printed in the problem and primary figure: a small-amplitude mechanical oscillator, the shown string route, 60.0 Hz, 3.5 × 10⁻⁴ kg/m, and the labelled 1.50 m horizontal span. No wavelength, speed, tension, loop count, or hanging-mass value occurs here
\end{definition}
\begin{proof}
  This is the declared model definition of Matches Problem And Primary Figure.
\end{proof}
\begin{definition}[Has Node Boundary Conditions]
  \label{def:physics:phyx-mini-0204:phyxminiproblems-problemphyxmini0204-hasnodeboundaryconditions}
  \lean{PhyXMiniProblems.ProblemPhyXMini0204.HasNodeBoundaryConditions}
  \uses{def:physics:phyx-mini-0204:phyxminiproblems-problemphyxmini0204-stringpulleysetup}
  Boundary-node idealization for the vibrating horizontal span. The oscillator node is explicitly stipulated in the source; the pulley contact constrains the other end of the idealized node-to-node span
\end{definition}
\begin{proof}
  This is the declared model definition of Has Node Boundary Conditions.
\end{proof}
\begin{definition}[Is Requested One Loop State]
  \label{def:physics:phyx-mini-0204:phyxminiproblems-problemphyxmini0204-isrequestedoneloopstate}
  \lean{PhyXMiniProblems.ProblemPhyXMini0204.IsRequestedOneLoopState}
  \uses{def:physics:phyx-mini-0204:phyxminiproblems-problemphyxmini0204-stringpulleysetup}
  The operating state requested by the question: exactly one loop
\end{definition}
\begin{proof}
  This is the declared model definition of Is Requested One Loop State.
\end{proof}
\begin{definition}[Uses Standard Gravity]
  \label{def:physics:phyx-mini-0204:phyxminiproblems-problemphyxmini0204-usesstandardgravity}
  \lean{PhyXMiniProblems.ProblemPhyXMini0204.UsesStandardGravity}
  \uses{def:physics:phyx-mini-0204:phyxminiproblems-problemphyxmini0204-accelerationinmeterspersecondsquared, def:physics:phyx-mini-0204:phyxminiproblems-problemphyxmini0204-stringpulleysetup}
  The standard classroom calibration g = 9.8 m/s². It is auxiliary physical data, kept separate from the values printed in the source and figure
\end{definition}
\begin{proof}
  This is the declared model definition of Uses Standard Gravity.
\end{proof}
\begin{definition}[Has Positive Physical Parameters]
  \label{def:physics:phyx-mini-0204:phyxminiproblems-problemphyxmini0204-haspositivephysicalparameters}
  \lean{PhyXMiniProblems.ProblemPhyXMini0204.HasPositivePhysicalParameters}
  \uses{def:physics:phyx-mini-0204:phyxminiproblems-problemphyxmini0204-lengthinmeters, def:physics:phyx-mini-0204:phyxminiproblems-problemphyxmini0204-massinkilograms, def:physics:phyx-mini-0204:phyxminiproblems-problemphyxmini0204-frequencyinhertz, def:physics:phyx-mini-0204:phyxminiproblems-problemphyxmini0204-linearmassdensityinkilogramspermeter, def:physics:phyx-mini-0204:phyxminiproblems-problemphyxmini0204-speedinmeterspersecond, def:physics:phyx-mini-0204:phyxminiproblems-problemphyxmini0204-forceinnewtons, def:physics:phyx-mini-0204:phyxminiproblems-problemphyxmini0204-accelerationinmeterspersecondsquared, def:physics:phyx-mini-0204:phyxminiproblems-problemphyxmini0204-stringpulleysetup}
  Positivity conditions for a taut string carrying a nonzero load
\end{definition}
\begin{proof}
  This is the declared model definition of Has Positive Physical Parameters.
\end{proof}
\begin{definition}[Satisfies Node To Node Standing Wave Law]
  \label{def:physics:phyx-mini-0204:phyxminiproblems-problemphyxmini0204-satisfiesnodetonodestandingwavelaw}
  \lean{PhyXMiniProblems.ProblemPhyXMini0204.SatisfiesNodeToNodeStandingWaveLaw}
  \uses{def:physics:phyx-mini-0204:phyxminiproblems-problemphyxmini0204-lengthinmeters, def:physics:phyx-mini-0204:phyxminiproblems-problemphyxmini0204-stringpulleysetup, def:physics:phyx-mini-0204:phyxminiproblems-problemphyxmini0204-hasnodeboundaryconditions}
  Node-to-node standing-wave geometry: a mode with n loops satisfies n λ = 2 ℓ. This is stated generically in the loop count rather than baking the requested one-loop answer into the law
\end{definition}
\begin{proof}
  This is the declared model definition of Satisfies Node To Node Standing Wave Law.
\end{proof}
\begin{definition}[Satisfies Wave Speed Law]
  \label{def:physics:phyx-mini-0204:phyxminiproblems-problemphyxmini0204-satisfieswavespeedlaw}
  \lean{PhyXMiniProblems.ProblemPhyXMini0204.SatisfiesWaveSpeedLaw}
  \uses{def:physics:phyx-mini-0204:phyxminiproblems-problemphyxmini0204-lengthinmeters, def:physics:phyx-mini-0204:phyxminiproblems-problemphyxmini0204-frequencyinhertz, def:physics:phyx-mini-0204:phyxminiproblems-problemphyxmini0204-speedinmeterspersecond, def:physics:phyx-mini-0204:phyxminiproblems-problemphyxmini0204-stringpulleysetup}
  The nondispersive wave relation v = f λ on the horizontal string
\end{definition}
\begin{proof}
  This is the declared model definition of Satisfies Wave Speed Law.
\end{proof}
\begin{definition}[Satisfies Taut String Tension Law]
  \label{def:physics:phyx-mini-0204:phyxminiproblems-problemphyxmini0204-satisfiestautstringtensionlaw}
  \lean{PhyXMiniProblems.ProblemPhyXMini0204.SatisfiesTautStringTensionLaw}
  \uses{def:physics:phyx-mini-0204:phyxminiproblems-problemphyxmini0204-linearmassdensityinkilogramspermeter, def:physics:phyx-mini-0204:phyxminiproblems-problemphyxmini0204-speedinmeterspersecond, def:physics:phyx-mini-0204:phyxminiproblems-problemphyxmini0204-forceinnewtons, def:physics:phyx-mini-0204:phyxminiproblems-problemphyxmini0204-stringpulleysetup}
  The transverse-wave law for a uniform taut string, written as T = μ v² in coherent SI readouts
\end{definition}
\begin{proof}
  This is the declared model definition of Satisfies Taut String Tension Law.
\end{proof}
\begin{definition}[Satisfies Ideal Pulley And Static Load Law]
  \label{def:physics:phyx-mini-0204:phyxminiproblems-problemphyxmini0204-satisfiesidealpulleyandstaticloadlaw}
  \lean{PhyXMiniProblems.ProblemPhyXMini0204.SatisfiesIdealPulleyAndStaticLoadLaw}
  \uses{def:physics:phyx-mini-0204:phyxminiproblems-problemphyxmini0204-massinkilograms, def:physics:phyx-mini-0204:phyxminiproblems-problemphyxmini0204-forceinnewtons, def:physics:phyx-mini-0204:phyxminiproblems-problemphyxmini0204-accelerationinmeterspersecondsquared, def:physics:phyx-mini-0204:phyxminiproblems-problemphyxmini0204-stringpulleysetup}
  Ideal frictionless-pulley transmission together with static force balance of the suspended mass: the horizontal and vertical tensions agree and the latter equals the weight m g
\end{definition}
\begin{proof}
  This is the declared model definition of Satisfies Ideal Pulley And Static Load Law.
\end{proof}
\begin{lemma}[hanging Mass In Kilograms eq exact]
  \label{lem:physics:phyx-mini-0204:phyxminiproblems-problemphyxmini0204-hangingmassinkilograms-eq-exact}
  \lean{PhyXMiniProblems.ProblemPhyXMini0204.hangingMassInKilograms_eq_exact}
  \uses{def:physics:phyx-mini-0204:phyxminiproblems-problemphyxmini0204-massinkilograms, def:physics:phyx-mini-0204:phyxminiproblems-problemphyxmini0204-stringpulleysetup, def:physics:phyx-mini-0204:phyxminiproblems-problemphyxmini0204-matchesproblemandprimaryfigure, def:physics:phyx-mini-0204:phyxminiproblems-problemphyxmini0204-hasnodeboundaryconditions, def:physics:phyx-mini-0204:phyxminiproblems-problemphyxmini0204-isrequestedoneloopstate, def:physics:phyx-mini-0204:phyxminiproblems-problemphyxmini0204-usesstandardgravity, def:physics:phyx-mini-0204:phyxminiproblems-problemphyxmini0204-haspositivephysicalparameters, def:physics:phyx-mini-0204:phyxminiproblems-problemphyxmini0204-satisfiesnodetonodestandingwavelaw, def:physics:phyx-mini-0204:phyxminiproblems-problemphyxmini0204-satisfieswavespeedlaw, def:physics:phyx-mini-0204:phyxminiproblems-problemphyxmini0204-satisfiestautstringtensionlaw, def:physics:phyx-mini-0204:phyxminiproblems-problemphyxmini0204-satisfiesidealpulleyandstaticloadlaw}
  The governing laws give the unrounded suspended mass 81/70 kg, approximately 1.157 kg, before comparison with the displayed one-decimal choices
\end{lemma}
\begin{proof}
  The conclusion follows from the governing relations and figure readouts named in the statement of hanging Mass In Kilograms eq exact.
\end{proof}
\begin{definition}[Answer Choice]
  \label{def:physics:phyx-mini-0204:phyxminiproblems-problemphyxmini0204-answerchoice}
  \lean{PhyXMiniProblems.ProblemPhyXMini0204.AnswerChoice}
  Labels of the four answers printed with the problem
\end{definition}
\begin{proof}
  This is the declared model definition of Answer Choice.
\end{proof}
\begin{definition}[kilograms]
  \label{def:physics:phyx-mini-0204:phyxminiproblems-problemphyxmini0204-answerchoice-kilograms}
  \lean{PhyXMiniProblems.ProblemPhyXMini0204.AnswerChoice.kilograms}
  \uses{def:physics:phyx-mini-0204:phyxminiproblems-problemphyxmini0204-answerchoice}
  Mass in kilograms printed beside each answer label
\end{definition}
\begin{proof}
  This is the declared model definition of kilograms.
\end{proof}
\begin{definition}[recorded Answer Choice]
  \label{def:physics:phyx-mini-0204:phyxminiproblems-problemphyxmini0204-recordedanswerchoice}
  \lean{PhyXMiniProblems.ProblemPhyXMini0204.recordedAnswerChoice}
  \uses{def:physics:phyx-mini-0204:phyxminiproblems-problemphyxmini0204-answerchoice}
  The answer label recorded in the supplied dataset metadata
\end{definition}
\begin{proof}
  This is the declared model definition of recorded Answer Choice.
\end{proof}
\begin{definition}[Matches Answer Choice]
  \label{def:physics:phyx-mini-0204:phyxminiproblems-problemphyxmini0204-matchesanswerchoice}
  \lean{PhyXMiniProblems.ProblemPhyXMini0204.MatchesAnswerChoice}
  \uses{def:physics:phyx-mini-0204:phyxminiproblems-problemphyxmini0204-massquantity, def:physics:phyx-mini-0204:phyxminiproblems-problemphyxmini0204-massinkilograms, def:physics:phyx-mini-0204:phyxminiproblems-problemphyxmini0204-answerchoice}
  Agreement with a displayed one-decimal kilogram answer. A 0.05 kg tolerance models rounding to the nearest tenth rather than identifying the exact physical mass with the printed decimal
\end{definition}
\begin{proof}
  This is the declared model definition of Matches Answer Choice.
\end{proof}
% --- Archon declaration topology end ---
% --- Archon physics formalization source end ---
